% !TEX root = ../../report.tex

\section{Анализ предметной области}

\textbf{Онлайн-маркетплейс} --- электронная торговая площадка, на которой множество независимых продавцов размещают свои товары, а покупатели выбирают и заказывают их через единый интерфейс~\cite{marketplace_def}. В отличие от классического интернет-магазина, маркетплейс не владеет товарами и не осуществляет их закупку: платформа выступает посредником между продавцом и покупателем, предоставляя инфраструктуру для размещения каталога, приёма заказов и сбора обратной связи.

По данным Ассоциации компаний интернет-торговли (АКИТ), объём интернет-торговли в России по итогам 2025 года составил 11,5~трлн рублей, что соответствует росту на 28\% относительно предыдущего года~\cite{akit2025}. Доля онлайн-продаж достигла 18,8\% от всего розничного товарооборота. Маркетплейсы являются значимой частью рынка онлайн-торговли.

Ключевые бизнес-процессы маркетплейса включают:

\begin{itemize}
    \item регистрацию продавцов и размещение товаров с привязкой к категориям;
    \item просмотр каталога покупателями с возможностью фильтрации и поиска;
    \item оформление заказов с фиксацией цены на момент покупки;
    \item управление жизненным циклом заказа (от корзины до доставки);
    \item систему отзывов и рейтингов, обеспечивающую обратную связь;
    \item формирование аналитической отчётности для продавцов и аналитиков.
\end{itemize}

Рассмотренные далее маркетплейсы являются закрытыми коммерческими системами. В проектируемой системе отдельно рассматриваются каталог, заказы, отзывы, история изменения цен и отчётность продавца.
